\section{Introduction}
\label{sec:introduction}

\subsection{Why Resolution Matters}

Congressional redistricting algorithms operate on graphs.
The vertices are geographic units; the edges are shared boundaries.
The choice of geographic unit --- the \emph{resolution} of the plan ---
determines the size of the graph, the granularity with which district
boundaries can be drawn, and the computational cost of every algorithm
that touches the graph.

The U.S. Census Bureau publishes three nested levels of geographic
unit that are practically useful for redistricting:

\begin{itemize}
  \item \textbf{Block groups} (\bg{}): approximately 600--3{,}000 persons,
        median $\approx$1{,}000.
        The finest unit in the P.L.\ 94-171 redistricting file that carries
        both population and race/ethnicity tabulations.
        North Carolina contains approximately 5{,}000 block groups.
  \item \textbf{Census tracts} (\tract{}): approximately 1{,}200--8{,}000 persons,
        median $\approx$4{,}000.
        The conventional redistricting unit; the current pipeline default.
        North Carolina contains 2{,}195 tracts.
  \item \textbf{Counties} (\county{}): typically 10{,}000--500{,}000 persons,
        state average $\approx$100{,}000.
        The coarsest useful level; 100 counties in North Carolina.
\end{itemize}

\noindent
Each level trades spatial precision against computational cost.
Block groups allow finer boundaries but approximately double the graph size
relative to tracts; counties provide a radical compression to roughly 5\%
of the tract graph size.
Neither extreme is universally correct.
State legislative chambers with many small districts (e.g., a 120-seat
state house) require finer resolution to prevent population equality
violations within a single tract from being unresolvable; large-$k$
congressional states (Texas, $k=38$) benefit from county-level coarsening
that reduces MCMC autocorrelation at modest precision cost.

\subsection{The GEOID Hierarchy as Free Information}

The Census Bureau encodes the nesting of geographic units directly in
the GEOID identifier~\citep{census2020geoids}.
Every block-group GEOID is a one-character extension of a tract GEOID;
every tract GEOID is a six-character extension of a county GEOID:

\begin{center}
\begin{tabular}{llll}
\toprule
Level & GEOID length & Example (NC, Wake County, tract 0531.01, BG 1) & Structure \\
\midrule
County  & 5 chars  & \geoid{37183}              & state(2)+county(3) \\
Tract   & 11 chars & \geoid{37183053101}        & state(2)+county(3)+tract(6) \\
Block group & 12 chars & \geoid{371830531011}   & state(2)+county(3)+tract(6)+bg(1) \\
\bottomrule
\end{tabular}
\end{center}

This prefix structure means that the complete fine-to-coarse mapping
at any adjacent pair of levels --- BG$\to$tract, tract$\to$county,
BG$\to$county --- is derivable from the GEOIDs already present in the
adjacency files.
No extra mapping files are required.
This observation is the key insight enabling the resolution system:
\emph{geographic hierarchy is encoded in the data we already have}.

\subsection{Three Problems Solved}

The current pipeline exposes a \texttt{--resolution block\_group} flag that
affects only adjacency loading; resolution does not flow through plan
generation, analysis, reporting, or audit~\citep{dellaLibera2026resolutionSpec}.
This creates three concrete problems:

\begin{enumerate}
  \item \textbf{Multi-scale MCMC requires two simultaneous resolutions}.
        The Autry et al.\ metropolized multiscale algorithm~\citep{autry2021multiscale}
        runs a fine-level ReCom chain with coarse-level acceptance/rejection
        guidance.
        The current pipeline has no concept of a resolution hierarchy ---
        it cannot represent ``fine = tract, coarse = county'' as a
        pipeline-level parameter.
  \item \textbf{High-precision redistricting cannot output BG-level plans}.
        A state house with 120 districts and average district population
        $\approx$86{,}000 cannot be resolved at tract level in dense urban
        counties; BG-level output is needed but not supported.
  \item \textbf{County-level coarsening has no clean integration path}.
        The county adjacency graph can be derived from the tract graph with
        zero extra data files, yet there is no mechanism to build it
        on-demand or to record the derivation formula in the audit chain.
\end{enumerate}

This paper addresses all three by making resolution a first-class parameter
from data loading through to the cryptographic audit manifest.

\subsection{Contributions}

The contributions of this paper are:

\begin{enumerate}
  \item \textbf{\derive{}}: a GEOID prefix-matching function that derives
        any fine-to-coarse partition, with a correctness proof and an
        explicit orphan detection guarantee (Section~\ref{sec:geoid}).
  \item \textbf{\buildc{}}: a function that builds the county adjacency graph
        from the tract graph with a formal proof of the adjacency criterion
        (Section~\ref{sec:geoid}).
  \item \textbf{Resolution threading}: \texttt{plan\_resolution}, \texttt{n\_units},
        and \texttt{unit\_type} fields in every plan manifest, enabling
        resolution-aware analysis and reporting (Section~\ref{sec:manifests}).
  \item \textbf{Three multi-scale options}: A (BG$\to$tract), B (tract$\to$county,
        implemented), C (BG$\to$county) — with data requirements, output levels,
        and audit formulas for each (Section~\ref{sec:options}).
  \item \textbf{Empirical validation}: Option~B reduces lag-100 autocorrelation
        on TX $k=38$; NC manifest examples for all three options
        (Section~\ref{sec:empirical}).
\end{enumerate}

\subsection{Paper Organization}

Section~\ref{sec:geoid} develops the GEOID partition derivation machinery:
formal definitions, correctness proofs, and the county coarsening criterion.
Section~\ref{sec:options} describes the three multi-scale options, their
CLI flags, and the implementation status of each.
Section~\ref{sec:manifests} describes the audit manifest extensions.
Section~\ref{sec:empirical} presents empirical results.
Section~\ref{sec:conclusion} concludes.
